% ============================================================================
% FORMATO - Tesis UTEM
% ============================================================================
% Configuración visual del documento: interlineado, paginación, estilos de
% títulos y encabezados, según la Pauta UTEM.
% ============================================================================

% --- Interlineado 1.5 líneas ---
\onehalfspacing

% --- Sangría de párrafo: 5 espacios (~1.25 cm) según norma UTEM (sección 4.6) ---
\setlength{\parindent}{1.25cm}

% --- Espacio entre párrafos: "Un espacio entre párrafos" (sección 4.6) ---
% Se usa un valor moderado para separar párrafos sin romper la lectura.
\setlength{\parskip}{6pt plus 2pt minus 1pt}

% ============================================================================
% ESTILO DE PÁGINA (encabezado y pie)
% ============================================================================
% Paginación: ángulo inferior derecho (sección 4.5)
\pagestyle{fancy}
\fancyhf{}                                   % Limpia encabezados y pies
\fancyfoot[R]{\thepage}                      % Número de página: ángulo inferior derecho
\renewcommand{\headrulewidth}{0pt}           % Sin línea en el encabezado
\renewcommand{\footrulewidth}{0pt}           % Sin línea en el pie

% Estilo para la primera página de cada capítulo
\fancypagestyle{plain}{
    \fancyhf{}
    \fancyfoot[R]{\thepage}
    \renewcommand{\headrulewidth}{0pt}
    \renewcommand{\footrulewidth}{0pt}
}

% ============================================================================
% ESTILO DE CAPÍTULOS
% ============================================================================
% Títulos de capítulos: centrados, mayúsculas, negrita, parte superior de la página
% Sección 4.6: "Inicio de los títulos en la parte superior de la página,
% en margen centrado, letras mayúsculas y en negrita."
\titleformat{\chapter}[display]
    {\normalfont\Large\bfseries\centering}    % Formato: grande, negrita, centrado
    {\MakeUppercase{\chaptertitlename}\ \thechapter} % "CAPÍTULO X"
    {10pt}                                    % Espacio entre número y título
    {\Large\MakeUppercase}                    % Título en mayúsculas

% Capítulos sin número (Introducción, Conclusiones, Bibliografía)
\titleformat{name=\chapter,numberless}[display]
    {\normalfont\Large\bfseries\centering}
    {}
    {0pt}
    {\Large\MakeUppercase}

% Espacio antes y después del título del capítulo
% Sección 4.6: "Inicio del texto después de dos espacios bajo el título"
% before-sep de -50pt compensa el espacio interno que LaTeX añade antes de los capítulos,
% logrando que el título inicie en el margen superior real (4 cm).
\titlespacing*{\chapter}{0pt}{-50pt}{20pt}

% ============================================================================
% ESTILO DE SECCIONES
% ============================================================================
% Sección 4.6: "Uso de minúsculas en subtítulos, con excepción de la primera
% letra de la primera palabra y sustantivos propios."
\titleformat{\section}
    {\normalfont\large\bfseries}             % Formato
    {\thesection}                            % Numeración (1.1, 1.2, ...)
    {1em}                                    % Espacio después del número
    {}                                       % Sin transformación (minúsculas naturales)

% Subsecciones (ej: 1.1.1 Subtítulo)
\titleformat{\subsection}
    {\normalfont\normalsize\bfseries}
    {\thesubsection}
    {1em}
    {}

% Subsubsecciones (ej: 1.1.1.1)
\titleformat{\subsubsection}
    {\normalfont\normalsize\bfseries}
    {\thesubsubsection}
    {1em}
    {}

% ============================================================================
% TABLAS DE CONTENIDO, FIGURAS Y TABLAS
% ============================================================================
% Profundidad de numeración (hasta subsubsección: 1.1.1.1)
\setcounter{secnumdepth}{3}
% Profundidad en la tabla de contenido
\setcounter{tocdepth}{3}

% Redefinición de nombres para cumplir con la norma UTEM (usando babel)
\addto\captionsspanish{
    \renewcommand{\contentsname}{TABLA DE CONTENIDO}
    \renewcommand{\listtablename}{ÍNDICE DE TABLAS}
    \renewcommand{\listfigurename}{ÍNDICE DE ILUSTRACIONES}
    \renewcommand{\bibname}{BIBLIOGRAFÍA}
    \renewcommand{\appendixname}{ANEXO}
}


% ============================================================================
% PIES DE FIGURAS Y TABLAS
% ============================================================================
% Formato: "Figura 1. Descripción" / "Tabla 1. Descripción"
\captionsetup[figure]{labelsep=period, justification=centering}
\captionsetup[table]{labelsep=period, justification=centering}

% ============================================================================
% NOTAS AL PIE
% ============================================================================
% Sección 4.3: notas se escriben a espacio simple
\renewcommand{\footnotesize}{\small}

% ============================================================================
% LISTAS
% ============================================================================
% Personalización de listas con viñetas y numeradas
\setlist[itemize]{topsep=0pt, partopsep=0pt, parsep=0pt, itemsep=2pt}
\setlist[enumerate]{topsep=0pt, partopsep=0pt, parsep=0pt, itemsep=2pt}
